% =============================================================
% Relativistic Hydromagnetic Equations (Ch IV, §36--40)
% Agent 16: Covariant MHD, Alfven waves, causality
% =============================================================

\chapter{RELATIVISTIC HYDROMAGNETIC STABILITY}
\label{ch:rel4}
\setcounter{section}{35}

\section{Relativistic Hydromagnetics}\label{sec:rel36}

The classical treatment of hydromagnetics developed in Chapter~IV rests on
the non-relativistic approximation: displacement currents are discarded,
charge densities enter only at order $u^{2}/c^{2}$, and the Alfv\'en speed
$V_{A}$ is unconstrained.  When the magnetic energy density becomes
comparable to the rest-mass energy density---as in neutron-star
magnetospheres, relativistic jets, or gamma-ray burst outflows---none of
these simplifications is admissible.  One must work with the \emph{full}
Maxwell equations coupled to relativistic fluid dynamics through a covariant
energy-momentum tensor.

The relativistic formulation brings two qualitative changes:
\begin{enumerate}
\item The inertia of the fluid is measured not by $\rho$ alone but by the
      \emph{relativistic enthalpy density} $\enthalpy = \edensity + p$,
      where $\edensity$ includes the rest-mass energy.  Magnetic stresses
      contribute to the effective inertia as well.
\item All characteristic speeds---Alfv\'en, fast and slow
      magnetosonic---are \emph{automatically sub-luminal},
      $v < c$, as a consequence of the covariant structure.
\end{enumerate}

Throughout this chapter we adopt the metric signature $(-,+,+,+)$, keep
$c$ explicit, and normalise the four-velocity by $u^{\mu}u_{\mu} = -c^{2}$.
The projection tensor onto the local rest frame is
$\Delta^{\mu\nu} = g^{\mu\nu} + u^{\mu}u^{\nu}/c^{2}$.


\section{The relativistic equations of magnetohydrodynamics}\label{sec:rel37}

\subsection*{(a) \emph{Covariant Maxwell equations}}

The electromagnetic field is encoded in the antisymmetric Faraday tensor
$F^{\mu\nu}$.  Maxwell's equations in covariant form are
\begin{equation}\label{eq:rel4-1}
  \nabla_{\nu} F^{\mu\nu} = -4\pi J^{\mu},
\end{equation}
\begin{equation}\label{eq:rel4-2}
  \nabla_{\nu}\, {}^{*}\!F^{\mu\nu} = 0,
\end{equation}
where ${}^{*}\!F^{\mu\nu} = \tfrac{1}{2}\epsilon^{\mu\nu\alpha\beta}
F_{\alpha\beta}$ is the Hodge dual of $F^{\mu\nu}$, $J^{\mu}$ is the
four-current, and $\epsilon^{\mu\nu\alpha\beta}$ is the Levi-Civita
tensor.  Equation~\eqref{eq:rel4-2} is equivalent to
$\nabla_{[\alpha}F_{\beta\gamma]}=0$ (the Bianchi identity for $F$).

No approximation of discarding displacement currents is made: both
equations retain their full content.

\subsection*{(b) \emph{Electromagnetic quantities in the fluid frame}}

Given a fluid element with four-velocity $u^{\mu}$, one defines the
electric and magnetic four-vectors measured in its rest frame:
\begin{equation}\label{eq:rel4-3}
  e^{\mu} = F^{\mu\nu}\,u_{\nu},
  \qquad
  b^{\mu} = {}^{*}\!F^{\mu\nu}\,u_{\nu}.
\end{equation}
Both are orthogonal to $u^{\mu}$: $e^{\mu}u_{\mu}=0$, $b^{\mu}u_{\mu}=0$.
In the local rest frame, $e^{\mu}=(0,\mathbf{E}')$ and
$b^{\mu}=(0,\mathbf{B}')$ reduce to the ordinary three-vectors.

The Faraday tensor can be reconstructed from $(e^{\mu},b^{\mu})$:
\begin{equation}\label{eq:rel4-4}
  F^{\mu\nu} = \frac{1}{c}\bigl(u^{\mu}e^{\nu}-u^{\nu}e^{\mu}\bigr)
    - \epsilon^{\mu\nu\alpha\beta}\,\frac{u_{\alpha}\,b_{\beta}}{c}.
\end{equation}

\subsection*{(c) \emph{Ideal MHD condition}}

The ideal MHD limit corresponds to infinite conductivity, which in
covariant language is the requirement that the electric field vanish in the
fluid frame:
\begin{equation}\label{eq:rel4-5}
  e^{\mu} = F^{\mu\nu}\,u_{\nu} = 0.
\end{equation}
This is the relativistic counterpart of $\mathbf{E}+\mu\mathbf{u}\times
\mathbf{H}=0$ in Chandrasekhar's equation~(4).  When
equation~\eqref{eq:rel4-5} holds, the Faraday tensor is entirely determined
by $b^{\mu}$ and $u^{\mu}$:
\begin{equation}\label{eq:rel4-6}
  F^{\mu\nu} = -\epsilon^{\mu\nu\alpha\beta}\,
    \frac{u_{\alpha}\,b_{\beta}}{c}.
\end{equation}

\subsection*{(d) \emph{The total energy-momentum tensor}}

The energy-momentum tensor of the coupled fluid--electromagnetic system is
\begin{equation}\label{eq:rel4-7}
  T^{\mu\nu} = T^{\mu\nu}_{\mathrm{fluid}} + T^{\mu\nu}_{\mathrm{EM}}.
\end{equation}
For a perfect fluid,
\begin{equation}\label{eq:rel4-8}
  T^{\mu\nu}_{\mathrm{fluid}} = (\edensity+p)\,
    \frac{u^{\mu}u^{\nu}}{c^{2}} + p\,g^{\mu\nu}.
\end{equation}
The electromagnetic energy-momentum tensor is
\begin{equation}\label{eq:rel4-9}
  T^{\mu\nu}_{\mathrm{EM}} = \frac{1}{4\pi}\!\left(
    F^{\mu\alpha}F^{\nu}{}_{\alpha}
    - \frac{1}{4}\,g^{\mu\nu}F^{\alpha\beta}F_{\alpha\beta}\right).
\end{equation}
Under the ideal MHD condition~\eqref{eq:rel4-5}, the electromagnetic part
can be rewritten entirely in terms of $b^{\mu}$ and $\bsq \equiv b^{\mu}b_{\mu}$:
\begin{equation}\label{eq:rel4-10}
  T^{\mu\nu}_{\mathrm{EM}} = \frac{\bsq}{4\pi}\!\left(
    \frac{u^{\mu}u^{\nu}}{c^{2}} + \frac{1}{2}\,g^{\mu\nu}\right)
    - \frac{b^{\mu}b^{\nu}}{4\pi}.
\end{equation}
Combining equations~\eqref{eq:rel4-8} and \eqref{eq:rel4-10}, the total
energy-momentum tensor of ideal relativistic MHD is
\begin{equation}\label{eq:rel4-11}
  \boxed{
  T^{\mu\nu} = \left(\edensity + p + \frac{\bsq}{4\pi}\right)
    \frac{u^{\mu}u^{\nu}}{c^{2}}
    + \left(p + \frac{\bsq}{8\pi}\right)g^{\mu\nu}
    - \frac{b^{\mu}b^{\nu}}{4\pi}.
  }
\end{equation}
This is the fundamental dynamical equation of relativistic MHD.  The
magnetic field contributes to the effective energy density through
$\bsq/8\pi$, to the isotropic pressure through $\bsq/8\pi$, and exerts
an anisotropic stress $-b^{\mu}b^{\nu}/4\pi$ (tension along the field
lines), exactly paralleling the non-relativistic decomposition of the
Lorentz force in equation~(8) of Chapter~IV.

\subsection*{(e) \emph{Conservation laws and equations of motion}}

The equations of motion follow from
\begin{equation}\label{eq:rel4-12}
  \nabla_{\nu}T^{\mu\nu} = 0 \qquad \text{(energy-momentum conservation)},
\end{equation}
together with baryon number conservation
\begin{equation}\label{eq:rel4-13}
  \nabla_{\mu}(\rdensity\, u^{\mu}) = 0,
\end{equation}
and the homogeneous Maxwell equation~\eqref{eq:rel4-2}, which under ideal
MHD becomes a transport equation for $b^{\mu}$ (see \S\ref{sec:rel38}).

Projecting equation~\eqref{eq:rel4-12} along and perpendicular to
$u^{\mu}$ yields, respectively, the relativistic energy equation and the
relativistic Euler equation:
\begin{equation}\label{eq:rel4-14}
  u_{\mu}\nabla_{\nu}T^{\mu\nu} = 0 \quad \Longrightarrow \quad
  \text{energy conservation along flow lines},
\end{equation}
\begin{equation}\label{eq:rel4-15}
  \Delta^{\mu}{}_{\alpha}\,\nabla_{\nu}T^{\alpha\nu} = 0
  \quad \Longrightarrow \quad
  \text{relativistic MHD Euler equation}.
\end{equation}
Writing out equation~\eqref{eq:rel4-15} explicitly:
\begin{equation}\label{eq:rel4-16}
  \left(\edensity + p + \frac{\bsq}{4\pi}\right)
    \frac{u^{\nu}\nabla_{\nu}u^{\mu}}{c^{2}}
  = -\Delta^{\mu\nu}\!\left(\nabla_{\nu}p
    + \nabla_{\nu}\frac{\bsq}{8\pi}\right)
    + \frac{b^{\nu}\nabla_{\nu}b^{\mu}}{4\pi}
    + \frac{b^{\mu}}{4\pi c^{2}}\,b_{\nu}\,
      u^{\alpha}\nabla_{\alpha}u^{\nu}.
\end{equation}
The left-hand side shows that the effective inertia is
$(\edensity+p+\bsq/4\pi)/c^{2}$, not $\rho$ alone.

\begin{relcorrection}
In the non-relativistic limit ($\edensity\to\rho c^{2}$, $p\ll\rho c^{2}$,
$\bsq\ll\rho c^{2}$), equation~\eqref{eq:rel4-16} reduces to
Chandrasekhar's equation~(10):
\[
  \rho\frac{du_{i}}{dt} = -\frac{\partial}{\partial x_{i}}
    \Bigl(\frac{p}{\rho}+\mu\frac{|\mathbf{H}|^{2}}{8\pi\rho}\Bigr)
    + \frac{\mu H_{j}}{4\pi\rho}\frac{\partial H_{i}}{\partial x_{j}}
    + \nu\nabla^{2}u_{i}.
\]
\end{relcorrection}


\section{Relativistic magnetic field evolution and flux freezing}
\label{sec:rel38}

\subsection*{(a) \emph{The relativistic induction equation}}

In non-relativistic MHD the magnetic field evolves according to
Chandrasekhar's equation~(13).  The covariant generalisation follows from
projecting the homogeneous Maxwell equation~\eqref{eq:rel4-2} with the
ideal MHD constraint~\eqref{eq:rel4-5}.  One obtains the \emph{relativistic
induction equation}:
\begin{equation}\label{eq:rel4-17}
  \boxed{
  u^{\nu}\nabla_{\nu}b^{\mu}
  - b^{\nu}\nabla_{\nu}u^{\mu}
  + b^{\mu}\,(\nabla_{\nu}u^{\nu})
  + \frac{u^{\mu}}{c^{2}}\,b_{\nu}\,
    (u^{\alpha}\nabla_{\alpha}u^{\nu}) = 0.
  }
\end{equation}
The first two terms are the covariant analogue of
$\partial\mathbf{H}/\partial t - \operatorname{curl}(\mathbf{u}\times
\mathbf{H})$; the third accounts for compressibility (cf.\ the
$H_{i}/\rho$ formulation, equation~(37) of Ch.~IV); and the fourth is a
purely relativistic correction that maintains the orthogonality
$b^{\mu}u_{\mu}=0$ under Lorentz boosts.

In the non-relativistic limit the fourth term is $\order{v^{2}/c^{2}}$
and equation~\eqref{eq:rel4-17} reduces precisely to Chandrasekhar's
equation~(15).

\subsection*{(b) \emph{Joule dissipation with finite conductivity}}

When the conductivity $\sigma$ is finite, the ideal MHD
condition~\eqref{eq:rel4-5} is replaced by the relativistic Ohm's law:
\begin{equation}\label{eq:rel4-18}
  J^{\mu} + \frac{1}{c^{2}}\,(J^{\nu}u_{\nu})\,u^{\mu}
  = \sigma\, F^{\mu\nu}\,u_{\nu}
  = \sigma\, e^{\mu}.
\end{equation}
The left-hand side is the spatial current in the fluid frame.  The Joule
dissipation rate per unit proper volume is
\begin{equation}\label{eq:rel4-19}
  \dot{q}_{\mathrm{Joule}} = \frac{e^{\mu}e_{\mu}}{\sigma}
  = \frac{1}{\sigma}\,F^{\mu\alpha}u_{\alpha}\,F_{\mu}{}^{\beta}u_{\beta}
  \geq 0.
\end{equation}
This is manifestly Lorentz invariant and non-negative, generalising the
$|\mathbf{J}|^{2}/\sigma$ of equation~(21).  In the absence of motions it
drives the monotonic decay of magnetic energy, exactly as in
\S\,38(a) of Chapter~IV.

When resistive effects are included, the induction
equation~\eqref{eq:rel4-17} acquires diffusive terms proportional to
$\eta = c^{2}/(4\pi\sigma)$ (the relativistic magnetic diffusivity).  The
decay time for a mode of spatial scale $L$ remains $\tau\sim L^{2}/\eta$,
in agreement with equation~(30) of Chapter~IV.

\subsection*{(c) \emph{Magnetic flux freezing: the relativistic Alfv\'en theorem}}

In ideal MHD ($\sigma\to\infty$), equation~\eqref{eq:rel4-17} implies a
relativistic generalisation of the Helmholtz--Kelvin--Alfv\'en theorem
(equation~(33) of Ch.~IV).

Define the magnetic flux through a two-surface $\mathcal{S}$ bounded by a
closed curve $\mathcal{C}$ co-moving with the fluid:
\begin{equation}\label{eq:rel4-20}
  \Phi_{\mathcal{S}} = \int_{\mathcal{S}} {}^{*}\!F_{\mu\nu}\,d\Sigma^{\mu\nu},
\end{equation}
where $d\Sigma^{\mu\nu}$ is the surface element.  Then, under ideal MHD,
\begin{equation}\label{eq:rel4-21}
  \boxed{
  \frac{d\Phi_{\mathcal{S}}}{d\tau} = 0.
  }
\end{equation}
\emph{The magnetic flux through any co-moving surface is conserved along
the flow.}  This is the covariant statement that magnetic field lines are
``frozen into'' the fluid, the exact relativistic analogue of
Chandrasekhar's theorem~(33).

The proof follows from the Lie-dragging of ${}^{*}\!F$ along $u^{\mu}$:
using the ideal MHD condition~\eqref{eq:rel4-5} and Maxwell's
equation~\eqref{eq:rel4-2}, one shows $\Lie_{u}\,{}^{*}\!F=0$, where
$\Lie_{u}$ denotes the Lie derivative along the flow.

For a compressible fluid the ratio $b^{\mu}/\rdensity$ satisfies a
transport equation analogous to that for $H_{i}/\rho$ in
equation~(37) of Chapter~IV:
\begin{equation}\label{eq:rel4-22}
  u^{\nu}\nabla_{\nu}\!\left(\frac{b^{\mu}}{\rdensity}\right)
  = \frac{b^{\nu}}{\rdensity}\,\nabla_{\nu}u^{\mu}
    - \frac{u^{\mu}}{c^{2}}\,\frac{b_{\nu}}{\rdensity}\,
      u^{\alpha}\nabla_{\alpha}u^{\nu}.
\end{equation}
The last term is the relativistic correction; it vanishes in the Newtonian
limit.


\section{Relativistic Alfv\'en waves}\label{sec:rel39}

\subsection*{(a) \emph{Linearised relativistic MHD}}

We consider an infinite, homogeneous medium at rest with a uniform magnetic
field $b^{\mu}_{0}=(0,\mathbf{B}_{0})$ in the fluid frame.  Let
$\edensity_{0}$, $p_{0}$ denote the equilibrium energy density and
pressure.  Small perturbations $\delta u^{\mu}$, $\delta b^{\mu}$,
$\delta\edensity$, $\delta p$ are imposed.

The linearised relativistic Euler equation~\eqref{eq:rel4-16} and induction
equation~\eqref{eq:rel4-17} decouple into three families of modes:
\emph{Alfv\'en} (shear), \emph{fast magnetosonic}, and \emph{slow
magnetosonic}.  We treat each in turn.

\subsection*{(b) \emph{Alfv\'en waves: dispersion relation}}

The Alfv\'en mode is an incompressible, transverse perturbation (no density
or pressure change) with $\delta u^{\mu}$ and $\delta b^{\mu}$
perpendicular to both $\mathbf{B}_{0}$ and the wave vector $\mathbf{k}$.
Seeking plane-wave solutions $\propto\exp[i(k_{\mu}x^{\mu})]$, the
linearised equations yield the dispersion relation
\begin{equation}\label{eq:rel4-23}
  \boxed{
  \omega^{2} = \vA^{2}\,k^{2}\cos^{2}\theta,
  }
\end{equation}
where $\theta$ is the angle between $\mathbf{k}$ and $\mathbf{B}_{0}$,
and the \emph{relativistic Alfv\'en speed} is
\begin{equation}\label{eq:rel4-24}
  \boxed{
  \vA = \frac{B_{0}/\sqrt{4\pi}}
    {\sqrt{(\edensity_{0}+p_{0})/c^{2} + B_{0}^{2}/(4\pi c^{2})}}
  = \frac{c\,\sqrt{\bsq_{0}/4\pi}}
    {\sqrt{\edensity_{0}+p_{0}+\bsq_{0}/4\pi}}.
  }
\end{equation}
Here $\bsq_{0}=B_{0}^{2}$ is the rest-frame magnetic field squared.

\begin{relcorrection}
\textbf{Comparison with classical result.}
In the non-relativistic limit $\edensity_{0}\to\rho c^{2}$,
$p_{0}\ll\rho c^{2}$, $B_{0}^{2}/4\pi\ll\rho c^{2}$, one recovers
Chandrasekhar's equation~(60):
\[
  V_{A} = \sqrt{\frac{\mu}{4\pi\rho}}\,H\cos\vartheta.
\]
The relativistic formula~\eqref{eq:rel4-24} differs in two ways:
(i) $\rho$ is replaced by $(\edensity_{0}+p_{0})/c^{2}$, reflecting
the full relativistic inertia, and (ii) the magnetic energy density itself
contributes to the inertia through the $B_{0}^{2}/(4\pi c^{2})$ term in
the denominator.
\end{relcorrection}

\subsection*{(c) \emph{Proof that $\vA < c$ always holds}}

From equation~\eqref{eq:rel4-24},
\begin{equation}\label{eq:rel4-25}
  \frac{\vA^{2}}{c^{2}}
  = \frac{\bsq_{0}/4\pi}{\edensity_{0}+p_{0}+\bsq_{0}/4\pi}.
\end{equation}
Since $\edensity_{0}+p_{0}>0$ (required by the dominant energy condition),
the denominator strictly exceeds the numerator, hence
\begin{equation}\label{eq:rel4-26}
  \vA^{2} < c^{2} \quad \text{for all physical configurations}.
\end{equation}
In the ultra-magnetised limit $\bsq_{0}/4\pi\gg\edensity_{0}+p_{0}$,
$\vA\to c$ from below but never reaches $c$.  This is a direct consequence
of the magnetic energy contributing to the inertia: the field's own energy
resists its own acceleration.

\begin{causalitycheck}
The relativistic Alfv\'en speed~\eqref{eq:rel4-24} is manifestly sub-luminal.
The phase speed of Alfv\'en waves,
$v_{\mathrm{ph}}=\omega/k=\vA|\cos\theta|\leq\vA<c$,
and the group speed
$|\partial\omega/\partial\mathbf{k}|=\vA<c$ are both causal.
\end{causalitycheck}

\subsection*{(d) \emph{Fast and slow magnetosonic waves}}

Including compressibility (density and pressure perturbations), two
additional families of modes emerge.  Define the sound speed
$\cs^{2}=(\partial p/\partial\edensity)_{s}\leq c^{2}$.
The fast ($+$) and slow ($-$) magnetosonic phase speeds for propagation at
angle $\theta$ to $\mathbf{B}_{0}$ satisfy
\begin{equation}\label{eq:rel4-27}
  v_{\pm}^{2} = \frac{1}{2}\!\left[
    \cs^{2}+\vA^{2}-\frac{\cs^{2}\vA^{2}}{c^{2}}
    \pm\sqrt{\left(\cs^{2}+\vA^{2}
      -\frac{\cs^{2}\vA^{2}}{c^{2}}\right)^{\!2}
      -4\,\cs^{2}\vA^{2}\cos^{2}\theta}
  \,\right].
\end{equation}
Setting $\theta=0$ (propagation along $\mathbf{B}_{0}$):
\[
  v_{+}^{2}\big|_{\theta=0}
  = \max(\cs^{2},\vA^{2}),
  \qquad
  v_{-}^{2}\big|_{\theta=0}
  = \min(\cs^{2},\vA^{2}).
\]

\begin{causalitycheck}
The fast magnetosonic speed satisfies
\begin{equation}\label{eq:rel4-28}
  \vf^{2} = \cs^{2} + \vA^{2} - \frac{\cs^{2}\vA^{2}}{c^{2}}
  = \cs^{2}+\vA^{2}\!\left(1-\frac{\cs^{2}}{c^{2}}\right) < c^{2},
\end{equation}
since $\cs<c$ and $\vA<c$.  The slow magnetosonic speed satisfies
$\vs^{2} = \cs^{2}\vA^{2}\cos^{2}\theta/\vf^{2} < \vf^{2} < c^{2}$.

\medskip\noindent
\textbf{All three MHD characteristic speeds---Alfv\'en, fast magnetosonic,
slow magnetosonic---are strictly sub-luminal.}  The system of relativistic
ideal MHD is therefore \emph{hyperbolic} and all signal propagation is
causal.
\end{causalitycheck}

\subsection*{(e) \emph{Effects of finite viscosity and resistivity}}

In the presence of shear viscosity $\shearvisc$ and resistivity $\eta$,
the Alfv\'en dispersion relation~\eqref{eq:rel4-23} generalises to
\begin{equation}\label{eq:rel4-29}
  (\omega - i\tilde{\nu}k^{2})(\omega - i\tilde{\eta}k^{2})
  = \vA^{2}k^{2}\cos^{2}\theta,
\end{equation}
where $\tilde{\nu}=\shearvisc/[(\edensity_{0}+p_{0}+\bsq_{0}/4\pi)/c^{2}]$
is the relativistic kinematic viscosity and $\tilde{\eta}=c^{2}/(4\pi\sigma)$
is the magnetic diffusivity.  This has exactly the structure of
Chandrasekhar's equation~(59), with the replacements $\rho\to
(\edensity_{0}+p_{0}+\bsq_{0}/4\pi)/c^{2}$ and $V_{A}\to\vA$.

The general solution is
\begin{equation}\label{eq:rel4-30}
  \omega = \pm k\!\left[\vA^{2}\cos^{2}\theta
    - \tfrac{1}{4}(\tilde{\nu}-\tilde{\eta})^{2}k^{2}\right]^{1/2}
  + \tfrac{1}{2}i(\tilde{\nu}+\tilde{\eta})k^{2}.
\end{equation}
Viscosity and resistivity damp the waves (additive effect) and modify the
effective phase speed (differential effect), precisely as in the classical
case (equation~(66) of Ch.~IV).

\begin{relcorrection}
\textbf{Israel--Stewart causal dissipation.}
The Navier--Stokes form of viscous and resistive damping used above is
acausal at high frequencies.  A fully consistent treatment replaces the
instantaneous constitutive relations with the Israel--Stewart relaxation
equations:
\[
  \taupi\, u^{\alpha}\nabla_{\alpha}\pi^{\langle\mu\nu\rangle}
    + \pi^{\mu\nu} = 2\shearvisc\,\sigma^{\mu\nu},
  \qquad
  \tau_{\eta}\,u^{\alpha}\nabla_{\alpha}e^{\langle\mu\rangle}
    + e^{\mu} = \sigma^{-1}\,J^{\langle\mu\rangle}_{\perp},
\]
where $\taupi>0$ and $\tau_{\eta}>0$ are relaxation times.  These render
the full dissipative MHD system hyperbolic, with all characteristic speeds
bounded by $c$.
\end{relcorrection}


\section{Equipartition, force-free fields, and the relativistic
Taylor--Proudman analogue}\label{sec:rel40}

\subsection*{(a) \emph{The relativistic equipartition solution}}

Just as in the classical case (\S\,40(a) of Ch.~IV), one may seek
solutions of the ideal relativistic MHD equations in which the velocity
field is everywhere aligned with the magnetic field.  Setting
\begin{equation}\label{eq:rel4-31}
  u^{\mu} = \pm\,\frac{c\,b^{\mu}}
    {\sqrt{(\edensity+p)\,c^{2}/(\bsq/4\pi)+\bsq}}\;
  = \pm\,\frac{b^{\mu}}
    {\sqrt{(\edensity+p)/(4\pi c^{2})+\bsq/(4\pi c^{2})}},
\end{equation}
one finds that the induction equation~\eqref{eq:rel4-17} is identically
satisfied and the Euler equation~\eqref{eq:rel4-16} reduces to a condition
of generalised hydrostatic equilibrium:
\begin{equation}\label{eq:rel4-32}
  \nabla_{\mu}\!\left(p + \frac{\bsq}{8\pi}\right) = 0.
\end{equation}
On this solution the kinetic and magnetic energy densities are in
equipartition:
\begin{equation}\label{eq:rel4-33}
  \frac{1}{2}\,\frac{(\edensity+p)}{c^{2}}\,|\mathbf{v}|^{2}
  = \frac{\bsq}{8\pi},
\end{equation}
which is the relativistic generalisation of Chandrasekhar's
equation~(73), with $\rho$ replaced by $(\edensity+p)/c^{2}$.

\begin{relcorrection}
In the non-relativistic limit the equipartition condition becomes
$\tfrac{1}{2}\rho|\mathbf{u}|^{2}=\mu|\mathbf{H}|^{2}/(8\pi)$,
recovering equation~(73) of Chapter~IV.
\end{relcorrection}

\subsection*{(b) \emph{Relativistic force-free fields}}

A force-free configuration requires $\nabla_{\nu}T^{\mu\nu}_{\mathrm{EM}}=0$,
i.e.\ the electromagnetic field exerts no net force on the fluid.  In
covariant language the condition is
\begin{equation}\label{eq:rel4-34}
  F^{\mu\nu}J_{\nu} = 0,
\end{equation}
meaning the four-current is null with respect to $F^{\mu\nu}$.  In the
fluid frame this requires $\mathbf{J}\times\mathbf{B}=\mathbf{0}$, so
that $\mathbf{J}\parallel\mathbf{B}$, exactly as in the classical
treatment (equation~(74) of Ch.~IV).

The simplest class of force-free solutions satisfies the relativistic
Beltrami condition:
\begin{equation}\label{eq:rel4-35}
  \nabla_{\nu}\,{}^{*}\!F^{\mu\nu} = 0,
  \qquad
  J^{\mu} = \alpha\, b^{\mu},
\end{equation}
where $\alpha$ is a scalar.  When $\alpha$ is constant this reduces to the
classical $\operatorname{curl}\mathbf{H}=\alpha\mathbf{H}$
(equation~(75)).

The virial constraint of equation~(79) carries over to the relativistic
case: a force-free field cannot vanish on the boundary of a finite volume
unless it vanishes everywhere.  The proof is identical, since it relies only
on the algebraic identity of the stress tensor, which has the same structure
in the relativistic case.

A generalisation permitting fluid motions is obtained by taking
$u^{\mu}=\beta\,c\,b^{\mu}/\sqrt{\bsq}$ with constant $\beta$, combined
with the force-free condition~\eqref{eq:rel4-34} and generalised
hydrostatic equilibrium, in exact analogy with equation~(80) of Chapter~IV.

\subsection*{(c) \emph{The relativistic analogue of the Taylor--Proudman
theorem}}

Consider a uniform background field $b^{\mu}_{0}$ and small steady
perturbations $\delta u^{\mu}$, $\delta b^{\mu}$.  In the ideal,
inviscid, non-resistive limit, the linearised Euler and induction equations
become
\begin{equation}\label{eq:rel4-36}
  \frac{B_{0,j}}{4\pi}\,\frac{\partial(\delta b_{i})}{\partial x_{j}}
  = \frac{\partial}{\partial x_{i}}
    \!\left(\frac{\delta p}{\edensity_{0}+p_{0}}
    +\frac{\mathbf{B}_{0}\cdot\delta\mathbf{b}}{4\pi(\edensity_{0}+p_{0})/c^{2}}\right),
\end{equation}
\begin{equation}\label{eq:rel4-37}
  B_{0,j}\,\frac{\partial(\delta u_{i})}{\partial x_{j}} = 0.
\end{equation}
Equation~\eqref{eq:rel4-37} states that \emph{steady, slow motions in the
presence of a uniform magnetic field cannot vary in the direction of
$\mathbf{B}_{0}$}.  This is the exact relativistic counterpart of
Chandrasekhar's equation~(83) and the hydromagnetic analogue of the
Taylor--Proudman theorem for rotating fluids.

\begin{relcorrection}
The Taylor--Proudman-type constraint~\eqref{eq:rel4-37} has the same
form in relativistic and classical MHD.  The relativistic modification
enters only through the pressure perturbation
in equation~\eqref{eq:rel4-36}, where $\rho$ is replaced by
$(\edensity_{0}+p_{0})/c^{2}$.  The two-dimensionality of the motions
is unaffected by relativistic corrections.
\end{relcorrection}


%% =============================================================
%% Summary of relativistic modifications, Chapter IV §36--40
%% =============================================================
\bigskip
\noindent\rule{\textwidth}{0.4pt}
\medskip
\noindent\textbf{Summary of key relativistic modifications
(Chapter IV, \S\S\,36--40).}

\begin{enumerate}
\item \textbf{Covariant Maxwell + MHD:} Full Maxwell equations are retained
      (no displacement-current approximation).  The electromagnetic field
      is encoded in $F^{\mu\nu}$; ideal MHD is $F^{\mu\nu}u_{\nu}=0$.
\item \textbf{Energy-momentum tensor:}
      $T^{\mu\nu}=(\edensity+p+\bsq/4\pi)\,u^{\mu}u^{\nu}/c^{2}
      +(p+\bsq/8\pi)\,g^{\mu\nu}-b^{\mu}b^{\nu}/4\pi$.
\item \textbf{Inertia:} $\rho \to (\edensity+p+\bsq/4\pi)/c^{2}$
      throughout; the magnetic field contributes to its own inertia.
\item \textbf{Alfv\'en speed:}
      $\vA = c\sqrt{\bsq_{0}/4\pi}/\sqrt{\edensity_{0}+p_{0}+\bsq_{0}/4\pi}
      < c$ always.
\item \textbf{Magnetosonic speeds:}
      $\vf^{2}=\cs^{2}+\vA^{2}-\cs^{2}\vA^{2}/c^{2}<c^{2}$;
      $\vs^{2}=\cs^{2}\vA^{2}\cos^{2}\theta/\vf^{2}<c^{2}$.
\item \textbf{Flux freezing:}
      $d\Phi/d\tau=0$ (relativistic Alfv\'en theorem).
\item \textbf{Dissipation:} Israel--Stewart causal relaxation replaces
      Navier--Stokes for finite $\nu$, $\eta$.
\item \textbf{Equipartition, force-free, Taylor--Proudman:} Classical
      results carry over with $\rho\to(\edensity+p)/c^{2}$.
\end{enumerate}
